\documentclass[a4paper,11pt]{article}
\usepackage[utf8]{inputenc}
\usepackage[T1]{fontenc}
\usepackage[french]{babel}
\usepackage[left=2cm,right=2cm,top=2cm,bottom=2cm]{geometry}
\usepackage{xcolor}
\usepackage{hyperref}
\usepackage{fontawesome5}
\usepackage{parskip}

% --- Couleurs ---
\definecolor{primary}{RGB}{35, 55, 123}

\begin{document}

% --- EN-TÊTE ---
\begin{center}
    {\Large \textbf{Loïc Aron MBASSI EWOLO}}\\[4pt]
    \small
    \faMapMarker\ Cergy, France \quad
    \faPhone\ (+33) 7 59 52 33 98 \quad
    \faEnvelope\ \href{mailto:loicaron.mbassiewolo@ensea.fr}{loicaron.mbassiewolo@ensea.fr}
\end{center}
\vspace{0.5cm}

% --- DESTINATAIRE ---
\hfill
\begin{minipage}[t]{0.45\textwidth}
    \textbf{ValoTec}\\
    Service Recrutement / R\&D\\
    Villejuif Bio Park\\
    1 mail du Professeur Georges Mathé\\
    94800 Villejuif
\end{minipage}

\vspace{1cm}

\textbf{Objet :} Candidature Stage Ingénieur R\&D Électronique (Avril 2026) + Projet Alternance\\
\textit{Réf : Offre Ingénieur R\&D et Industrialisation Électronique}

\vspace{0.5cm}

Madame, Monsieur,

L'intersection entre l'ingénierie électronique et le domaine médical est le fil conducteur que j'ai choisi pour ma carrière. Votre positionnement au sein du \textit{Villejuif Bio Park} et votre expertise dans le développement de dispositifs médicaux actifs font de ValoTec l'environnement idéal pour concrétiser cette ambition. Actuellement en double diplôme ingénieur à l'\textbf{ENSEA} (Cergy) et à l'\textbf{École Polytechnique de Yaoundé}, je souhaite mettre ma rigueur scientifique et mon savoir-faire technique au service de vos projets R\&D.

Mon profil se distingue par une approche concrète de l'ingénierie. Au laboratoire IOTA, j'ai contribué au projet \textit{Harmony Gloves} (gant connecté pour la langue des signes). Au sein de l'équipe, j'ai pris part activement à la conception matérielle, réalisant notamment la \textbf{soudure} des composants, ainsi qu'à la campagne de \textbf{prise de données} nécessaire à l'entraînement du modèle d'intelligence artificielle.

En parallèle de cette compétence hardware, je cultive un bagage algorithmique solide appliqué aux systèmes embarqués. Participant actuellement au challenge robotique \textit{CoHoMa} (Défense), j'y implémente des algorithmes de SLAM et des méthodes de Monte Carlo sur cibles Nvidia Jetson, en utilisant Docker pour garantir la fiabilité du déploiement en environnement contraint.

Cette approche terrain est complétée par une rigueur scientifique acquise lors de mon passage au \textbf{Mila (Institut Québécois d'IA)} en tant qu'\textbf{Assistant de Recherche}, où j'ai travaillé sur la généralisation des réseaux de neurones.

Souhaitant m'établir et poursuivre ma carrière d'ingénieur en France, je suis disponible dès \textbf{avril 2026} pour un stage \textbf{Assistant Ingénieur} de 4 à 5 mois. Mon objectif est de m'investir pleinement dans vos projets pour prouver ma valeur et poursuivre notre collaboration via un \textbf{Contrat de Professionnalisation} pour ma dernière année à l'ENSEA (dès septembre 2026).

Je serais ravi de vous présenter mes prototypes et de discuter de vos enjeux d'industrialisation lors d'un entretien.

Je vous prie d'agréer, Madame, Monsieur, l'expression de mes salutations distinguées.

\vspace{1cm}

\textbf{Loïc Aron MBASSI EWOLO}\\
\textit{Élève-Ingénieur ENSEA / Polytechnique Yaoundé}

\end{document}
